\clearpage\fancyhead[L]{\small PALLAS. COMET.}\hypertarget{native-332-lower}{}\label{native-332-lower}
\begin{center}\textbf{PALLAS. COMET.}\par\vspace{4pt}\hrule\vspace{4pt}\fontsize{10}{12}\selectfont Göttingische gelehrte Anzeigen. Issue 130. pp. 1289\,..\,1293. 17 August 1811.\par\vspace{4pt}\hrule\vspace{4pt}\end{center}

\par\indent Issue 8 of our Notices this year [p. 322 of this volume] reported Professor \textsc{Gauss}’s work on the perturbations of Pallas since its discovery. It also gave the \SourceLeaf{333}first observation, made on 4 January at our observatory, which confirmed the results of that work with a precision we could scarcely have expected. The planet was too faint for meridian observations at our observatory, so observations here went no further. The only step taken was to notify astronomers elsewhere in good time, since they had more powerful instruments at their disposal. For the period of the opposition, which occurred in February, Professor \textsc{Gauss} thus received very valuable observations from five observatories, in Paris, Mannheim, Seeberg, Berlin and Hamburg. These formed the basis of his further investigations. The observations made at the Seeberg and Mannheim observatories by Messrs. \textsc{von Lindenau}, \textsc{Barry} and \textsc{Harding} have already appeared in the Monatliche Correspondenz. We give the rest here.

\Needspace{14\baselineskip}\begin{center}\fontsize{10}{12}\selectfont\setlength{\tabcolsep}{4pt}\renewcommand{\arraystretch}{1.15}
Observations of Pallas in Paris, by Mr. \textsc{Bouvard}.

\begin{tabular}{c|c|c}
on the meridian & Right ascension & Southern declination \\
\hline
1811 Feb. 17 & $146^\circ\,58'\,33''15$ & $13^\circ\,29'\,42''4$ \\
18 & $146\;\;47\;\;55.95$ & $13\;\;\;5\;\;59.4$ \\
22 & $146\;\;\;6\;\;27.90$ & $11\;\;28\;\;19.6$ \\
27 & $145\;\;18\;\;13.65$ & $9\;\;19\;\;32.9$ \\
\end{tabular}
\end{center}

\Needspace{14\baselineskip}\begin{center}\fontsize{10}{12}\selectfont\setlength{\tabcolsep}{4pt}\renewcommand{\arraystretch}{1.15}
Observations in Berlin, by Professor \textsc{Bode}.

\begin{tabular}{c|c|c|c}
1811 & Mean time & Right ascension & Southern declination \\
\hline
Feb. 18 & $11^{\mathrm u}\,55^{\mathrm m}\,22^{\mathrm s}5$ & $146^\circ\,48'\,2''$ & --- --- --- \\
19 & $11\;\;50\;\;34.0$ & $146\;\;37\;\;32$ & $12^\circ\,42'\,44''5$ \\
20 & $11\;\;45\;\;57.0$ & $146\;\;27\;\;\;4$ & $12\;\;19\;\;\;1$ \\
March 12 & $10\;\;16\;\;33.6$ & $143\;\;45\;\;29$ & $3\;\;32\;\;42$ \\
16 & $9\;\;59\;\;45.1$ & $143\;\;29\;\;18$ & $1\;\;48\;\;\;3$ \\
17 & $9\;\;55\;\;35.1$ & $143\;\;26\;\;13$ & $1\;\;22\;\;11$ \\
\end{tabular}
\end{center}

\Needspace{14\baselineskip}\begin{center}\fontsize{10}{12}\selectfont\setlength{\tabcolsep}{4pt}\renewcommand{\arraystretch}{1.15}
Observations in Hamburg, by Professor \textsc{Schumacher}.

\begin{tabular}{c|c|c}
on the meridian & Right ascension & Southern declination \\
\hline
1811 Feb. 19 & $146^\circ\,37'\,15''0$ & $12^\circ\,42'\,45''3$ \\
20 & $146\;\;26\;\;52.5$ & $12\;\;18\;\;32.7$ \\
21 & $146\;\;16\;\;37.5$ & $11\;\;53\;\;47.7$ \\
22 & $146\;\;\;6\;\;36.0$ & $11\;\;28\;\;53.5$ \\
March 9 & $144\;\;\;1\;\;45.0$ & $4\;\;52\;\;42.5$ \\
15 & $143\;\;32\;\;39.0$ & $2\;\;13\;\;59.0$ \\
16 & $143\;\;29\;\;24.0$ & $1\;\;47\;\;57.0$ \\
18 & $143\;\;23\;\;34.5$ & $0\;\;56\;\;47.8$ \\
\end{tabular}
\end{center}


\Needspace{8\baselineskip}\noindent From these observations Professor \textsc{Gauss} derived the following result for the opposition.
\begin{center}1811 Feb. 21. $19^{\mathrm u}\,25^{\mathrm m}\,31^{\mathrm s}$ mean time in Göttingen\end{center}

\begin{center}\fontsize{10}{12}\selectfont\setlength{\tabcolsep}{4pt}\renewcommand{\arraystretch}{1.15}\begin{tabular}{llr}
True heliocentric longitude & .\;.\;.\;.\;.\;.\;.\;. & $152^\circ\,48'\,15''8$\\
True geocentric latitude, south & .\;.\;.\;.\;. & $23\;48\;19.2$\\
\end{tabular}\end{center}

\Needspace{17\baselineskip}\par\indent The agreement of this opposition with the previous six, once the perturbations caused by Jupiter had been duly applied as Professor \textsc{Gauss} had worked them out, is truly admirable. After only a few extremely small corrections had been made to the elements determined before the planet was found again, the differences between observation and calculation were as follows.
\SourceLeaf{334}
\begin{center}\fontsize{10}{12}\selectfont\setlength{\tabcolsep}{4pt}\renewcommand{\arraystretch}{1.15}
Difference in

\begin{tabular}{c|c|c}
 & mean longitude & heliocentric latitude \\
\hline
1803 & $+8''7$ & $-1''1$ \\
1804 & $-7.0$ & $-10.2$ \\
1805 & $+1.4$ & $-3.5$ \\
1807 & $-1.2$ & $-2.4$ \\
1808 & $-3.1$ & $-21.4$ \\
1809 & $+4.5$ & $-2.1$ \\
1811 & $-2.2$ & $-12.1$ \\
\end{tabular}
\end{center}


\Needspace{21\baselineskip}\par\indent The calculation of the perturbations was then carried forward another year. Professor \textsc{Gauss} was helped in this laborious work by Mr. \textsc{Nicolai}, who applies himself here with great zeal to the mathematical sciences. Under Professor \textsc{Gauss}’s supervision he carried out the greater part of the numerical calculations with as much diligence as skill. According to the results of these calculations, the motion of Pallas around the time of next year’s opposition can be represented by the following elliptic elements, which already include the perturbations.

\begin{center}\fontsize{10}{12}\selectfont\setlength{\tabcolsep}{4pt}\renewcommand{\arraystretch}{1.15}\begin{tabular}{llr}
\shortstack[l]{Epoch of mean longitude, 10 June 1812, noon\\at the Göttingen meridian} & . & $239^\circ\,4'\,46''1$\\
Longitude of perihelion for the same time &  & $121\;0\;48.5$\\
Longitude of the ascending node & . & $172\;32\;44.5$\\
Inclination of the orbit & .\;.\;.\;.\;.\;.\;. & $34^\circ\,34'\,54''7$\\
Mean daily tropical motion &  & $768''5746$\\
Eccentricity & .\;.\;.\;.\;.\;. & $=\sin13^\circ\,59'\,1''8$\\
Logarithm of the semimajor axis & . & $0.4429321$\\
\end{tabular}\end{center}

\Needspace{20\baselineskip}\par\indent Mr. \textsc{Nicolai} has also used these same elements to calculate an ephemeris for the apparition of Pallas in 1812. It is too long to include here and will be printed in the Monatl. Correspondenz and the Astronom. Jahrbuch. We give only the next opposition (the eighth), which Mr. \textsc{Nicolai} also calculated \emph{in advance}. In due course it can serve to test the accuracy of our results.
\begin{center}1812 June 10. $3^{\mathrm u}\,32^{\mathrm m}$ mean time in Göttingen.\end{center}

\begin{center}\fontsize{10}{12}\selectfont\setlength{\tabcolsep}{4pt}\renewcommand{\arraystretch}{1.15}\begin{tabular}{llr}
Heliocentric longitude of Pallas & .\;.\;.\;.\;.\;.\;.\;. & $259^\circ\,28'\,53''$\\
Heliocentric latitude, north & .\;.\;.\;.\;.\;.\;.\;. & $34\;32\;37$\\
Geocentric latitude & .\;.\;.\;.\;.\;.\;.\;.\;.\;.\;.\;. & $48\;16\;36$\\
\end{tabular}\end{center}

\begin{center}\fontsize{10}{12}\selectfont\setlength{\tabcolsep}{4pt}\renewcommand{\arraystretch}{1.15}\begin{tabular}{lr}
\shortstack[l]{Logarithm of the distance of Pallas\\from the Sun} & $0.504259$\\
\shortstack[l]{Logarithm of the distance of Pallas\\from Earth} & $0.384915$\\
\end{tabular}\end{center}

\Needspace{20\baselineskip}\par\indent We take this opportunity to add the parabolic elements of the comet discovered by \textsc{Flaugergues} in March of this year. Professor \textsc{Gauss} has just calculated them provisionally from some observations Mr. von \textsc{Zach} made in Marseille. They may make it easier to find this remarkable comet again toward the end of August.
\begin{center}Perihelion passage 1811 Sept. 10. $0^{\mathrm u}\,51^{\mathrm m}$\end{center}

\begin{center}\fontsize{10}{12}\selectfont\setlength{\tabcolsep}{4pt}\renewcommand{\arraystretch}{1.15}\begin{tabular}{llr}
Longitude of perihelion & .\;.\;.\;.\;.\;.\;.\;.\;.\;.\;.\;. & $73^\circ\,14'\,35''$\\
Logarithm of the least distance & .\;.\;.\;.\;.\;.\;. & $9.99153$\\
Longitude of the ascending node & .\;.\;.\;.\;.\;.\;.\;. & $141^\circ\,4'\,59''$\\
Inclination of the orbit & .\;.\;.\;.\;.\;.\;.\;.\;.\;.\;.\;.\;. & $73\;48\;2$\\
\emph{Retrograde} motion &  & \\
\end{tabular}\end{center}

\SourceLeaf{335}\Needspace{12\baselineskip}\par\indent The following positions were calculated from these elements.
\begin{center}\fontsize{10}{12}\selectfont\setlength{\tabcolsep}{4pt}\renewcommand{\arraystretch}{1.15}
\begin{tabular}{l|cc|c}
Midnight & Right ascension & Northern declination & Brightness \\
\hline
1811 Aug. 11 & $142^\circ\,43'$ & $29^\circ\,54'$ & $0.194$ \\
\phantom{1811 }\phantom{Aug. }26 & $152\;\;19$ & $36\;\;15$ & $0.294$ \\
Sept. 10 & $167\;\;34$ & $43\;\;14$ & $0.434$ \\
\phantom{Sept. }25 & $194\;\;22$ & $48\;\;35$ & $0.576$ \\
Oct. 10 & $232\;\;28$ & $44\;\;52$ & $0.595$ \\
\end{tabular}
\end{center}


\par\indent On the same scale the brightness (assuming that the comet receives its light solely from the Sun) was $0.035$ on 11 April and $0.052$ on 2 June. At the beginning of October it will therefore be 17 times as bright as on 11 April and will appear with exceptional brilliance in the northern sky. It will then pass through the constellations Hercules and Aquila as its light fades, and it will become invisible in Aquarius at the beginning of next year. If its steadily fading light has by then remained strong enough, it may become visible once more in Aquarius in April and may still be followed for quite some time with very powerful telescopes. Even if this proves impossible, its outstanding brilliance this autumn and its long period of visibility, which will surely amount to nearly ten months, will make it one of the most remarkable phenomena in the annals of cometography.
